\subsection*{Фаза 1. Сборка основы и ранняя интеграция}
\textbf{Период:} 2026-01-05 -- 2026-01-09 \\
\textbf{Коммитов в фазе:} 53

Проект собирается из нескольких заготовок в единую систему: складываются backend, UI, потоковый режим, \texttt{UMS} и первые пользовательские сценарии. По коммитам доминируют общая стабилизация и организационные правки и модели, \texttt{UMS}, распределение ресурсов и наблюдаемость. Это был период шлифовки, когда устойчивость становилась важнее простого прироста функций.

\begin{itemize}
\item \texttt{2b2010c} (2026-01-05) --- merge with all proj
\item \texttt{d0c788d} (2026-01-07) --- feat: Integrate Open WebUI, fix UMS inference, add MD reporting
\item \texttt{e7e7645} (2026-01-06) --- Chore: Clean up warnings and modernize FastAPI events in UMS
\item \texttt{f79bddb} (2026-01-06) --- Feature: Add .env support and clean up GPU monitoring imports
\item \texttt{62ff5e8} (2026-01-06) --- Feature: End-to-end real-time streaming
\end{itemize}

\subsection*{Фаза 2. Переход к v3.0 и аппаратно-адаптивной архитектуре}
\textbf{Период:} 2026-02-24 -- 2026-03-05 \\
\textbf{Коммитов в фазе:} 33

Появляется более зрелый контур \texttt{v3.0}: \texttt{Chainlit}, \texttt{AdaptiveRAGPipeline}, аппаратное профилирование, новые workflow и тестовое покрытие. По коммитам доминируют поддержка и эволюция \texttt{Chainlit}-контура и общая стабилизация и организационные правки. Это был период ускоренного наращивания продукта и архитектуры.

\begin{itemize}
\item \texttt{ceae15c} (2026-02-26) --- feat: v3.0 Hardware-Adaptive Agentic System — profiler, tiered RAG, Chainlit UI
\item \texttt{adf9bbf} (2026-03-04) --- chore(repo): ignore local test logs and ad-hoc node bootstrap
\item \texttt{512fd97} (2026-03-04) --- chore(repo): remove obsolete verification script
\item \texttt{0f59b70} (2026-02-26) --- chore: update .gitignore and requirements.txt for v3.0
\item \texttt{5c62f9a} (2026-03-04) --- docs(repo): consolidate task docs and ignore local artifacts
\end{itemize}

\subsection*{Фаза 3. Backend-first усиление и операционализация}
\textbf{Период:} 2026-03-12 -- 2026-03-31 \\
\textbf{Коммитов в фазе:} 124

Команда переводит проект от набора функций к платформе: усиливаются \texttt{orchestrator}, \texttt{UMS}, реестр моделей, \texttt{operator UI}, офлайн-поставка и наблюдаемость. По коммитам доминируют рантайм, запуск, офлайн-развёртывание и инсталляция и документация и фиксация правил эксплуатации. Это был период ускоренного наращивания продукта и архитектуры.

\begin{itemize}
\item \texttt{85bd9cb} (2026-03-18) --- feat(model-registry): add canonical yaml registry
\item \texttt{fac57e5} (2026-03-29) --- feat(operator): add python control plane for operator ui
\item \texttt{5ac5f34} (2026-03-31) --- Fix compare workflow LLM queue discrepancy and JSON parsing
\item \texttt{8ef99cd} (2026-03-17) --- chore(config): add local hardware override defaults
\item \texttt{d91b8f1} (2026-03-16) --- chore(git): ignore local runtime artifacts
\end{itemize}

\subsection*{Фаза 4. Консолидация вокруг Open WebUI-first}
\textbf{Период:} 2026-04-01 -- 2026-04-12 \\
\textbf{Коммитов в фазе:} 44

Основной вектор смещается к \texttt{Open WebUI}: уточняется разделение ответственности между инструментами, bootstrap-процессом, рантаймом и миграционной документацией. По коммитам доминируют документация и фиксация правил эксплуатации и миграция и стабилизация \texttt{Open WebUI}. Это был период ускоренного наращивания продукта и архитектуры.

\begin{itemize}
\item \texttt{71c9c51} (2026-04-11) --- feat(openwebui): reconcile bootstrap drift and native function-calling defaults
\item \texttt{4a800b1} (2026-04-01) --- Fix model download scripts: Qwen3-14B from MaziyarPanahi, LaBSE from cointegrated
\item \texttt{c8e8671} (2026-04-01) --- chore: update package version lower bounds and build script version pins
\item \texttt{062d6e9} (2026-04-07) --- docs(repo): relocate forms and research artifacts
\item \texttt{02d83c8} (2026-04-01) --- feat(compare): implement B3.51g and B3.51h - summary-first bounded latency with fast/deep modes
\end{itemize}
